
\documentclass[11pt,a4paper]{article}

\usepackage[utf8]{inputenc}
\usepackage[T1]{fontenc}
\usepackage[french]{babel}
\usepackage{lmodern}
\usepackage{microtype}
\usepackage[a4paper,margin=1.65cm,headheight=15pt]{geometry}
\usepackage{amsmath,amssymb,mathtools}
\usepackage{array,tabularx,multicol,booktabs}
\usepackage{enumitem}
\usepackage[most]{tcolorbox}
\usepackage{xcolor}
\usepackage{tikz}
\usetikzlibrary{arrows.meta,calc,positioning}
\usepackage{pdflscape}
\usepackage{fancyhdr}
\usepackage{hyperref}
\usepackage{needspace}

\definecolor{fondpage}{RGB}{247,248,250}
\definecolor{encre}{RGB}{31,35,40}
\definecolor{bleu}{RGB}{40,91,135}
\definecolor{vert}{RGB}{55,125,92}
\definecolor{orange}{RGB}{178,105,42}
\definecolor{violet}{RGB}{101,77,145}
\definecolor{rouge}{RGB}{170,72,72}
\definecolor{gris}{RGB}{86,91,99}
\definecolor{bleufonce}{RGB}{28,55,120}
\definecolor{bleuclair}{RGB}{238,243,255}
\definecolor{grisclair}{RGB}{245,245,245}
\definecolor{tableline}{RGB}{45,45,45}
\definecolor{softgray}{RGB}{246,247,249}
\definecolor{deepblue}{RGB}{25,62,115}
\definecolor{accent}{RGB}{20,82,130}

\hypersetup{colorlinks=true,linkcolor=bleufonce,urlcolor=bleufonce,
pdftitle={Parcours de revision -- Probabilités conditionnelles -- Premiere specialite},pdfauthor={}}

\setlength{\parindent}{0pt}
\setlength{\parskip}{0.25em}
\renewcommand{\arraystretch}{1.13}
\setlength{\tabcolsep}{4pt}
\setlist[itemize]{leftmargin=1.25em,itemsep=0.15em,topsep=0.2em}
\setlist[enumerate,1]{label=\textbf{\arabic*.}, leftmargin=1.05cm, itemsep=0.35em, topsep=0.2em}
\setlist[enumerate,2]{label=\textbf{\alph*.}, leftmargin=0.85cm, itemsep=0.25em, topsep=0.15em}

\pagestyle{fancy}
\fancyhf{}
\cfoot{\thepage}

\newcommand{\R}{\mathbb{R}}
\newcommand{\N}{\mathbb{N}}
\newcommand{\sourceq}[1]{ {\scriptsize\textcolor{bleufonce}{(site Première q#1)}}}

\newcounter{question}
\newcounter{reponse}

\newenvironment{blocquestions}[1]{%
  \begin{tcolorbox}[enhanced,breakable,colback=bleuclair,colframe=bleufonce,boxrule=0.6pt,arc=2mm,left=2mm,right=2mm,top=2mm,bottom=1.5mm,title={\bfseries #1},coltitle=white,colbacktitle=bleufonce,before skip=3mm,after skip=3mm, fontupper=\large]
  \large
  \begin{enumerate}[leftmargin=*,label=\textbf{\arabic*.},itemsep=0.85em,topsep=0.5em]
  \setcounter{enumi}{\value{question}}
}{%
  \setcounter{question}{\value{enumi}}
  \end{enumerate}
  \end{tcolorbox}
}

\newenvironment{blocreponses}[1]{%
  \begin{tcolorbox}[enhanced,breakable,colback=bleuclair,colframe=bleufonce,boxrule=0.6pt,arc=2mm,left=2mm,right=2mm,top=2mm,bottom=1.5mm,title={\bfseries #1},coltitle=white,colbacktitle=bleufonce,before skip=3mm,after skip=3mm, fontupper=\large]
  \large
  \begin{enumerate}[leftmargin=*,label=\textbf{\arabic*.},itemsep=0.45em,topsep=0.5em]
  \setcounter{enumi}{\value{reponse}}
}{%
  \setcounter{reponse}{\value{enumi}}
  \end{enumerate}
  \end{tcolorbox}
}

\newtcolorbox{correctionbox}[1]{enhanced,breakable,colback=grisclair,colframe=black!55,boxrule=0.55pt,arc=2mm,left=2mm,right=2mm,top=2mm,bottom=1.5mm,title=\textbf{#1},coltitle=black,colbacktitle=black!12,before skip=2mm,after skip=2mm}

\newtcolorbox{methodebox}[1]{enhanced,breakable,colback=white,colframe=bleu!70!black,boxrule=0.55pt,arc=2mm,left=2mm,right=2mm,top=2mm,bottom=1.5mm,title=\textbf{#1},coltitle=white,colbacktitle=bleu!70!black,before skip=2mm,after skip=2mm}

\newcommand{\titreSujet}[2]{%
  \subsection*{#1}
  \addcontentsline{toc}{subsection}{#1}
  \textit{Référence / inspiration : #2}\par\medskip\hrule\medskip
}

\newcommand{\qcm}[4]{%
\begin{center}
\begin{tabularx}{0.96\linewidth}{@{}XX@{}}
\textbf{a.} #1 & \textbf{b.} #2\\[1mm]
\textbf{c.} #3 & \textbf{d.} #4
\end{tabularx}
\end{center}}

\newcommand{\arbreDeuxNiveaux}[6]{%
\begin{center}
\begin{tikzpicture}[x=1cm,y=1cm,>=Latex]
\coordinate (O) at (0,0);
\coordinate (A) at (2,1.25);
\coordinate (Ab) at (2,-1.25);
\coordinate (AB) at (4.3,1.85);
\coordinate (ABb) at (4.3,0.65);
\coordinate (AbB) at (4.3,-0.65);
\coordinate (AbBb) at (4.3,-1.85);
\draw[->,line width=.55pt] (O)--(A) node[midway,above left,scale=.86] {#1};
\draw[->,line width=.55pt] (O)--(Ab) node[midway,below left,scale=.86] {#2};
\draw[->,line width=.55pt] (A)--(AB) node[midway,above,scale=.86] {#3};
\draw[->,line width=.55pt] (A)--(ABb) node[midway,below,scale=.86] {#4};
\draw[->,line width=.55pt] (Ab)--(AbB) node[midway,above,scale=.86] {#5};
\draw[->,line width=.55pt] (Ab)--(AbBb) node[midway,below,scale=.86] {#6};
\node[font=\small\bfseries] at (A) [right=1mm] {$A$};
\node[font=\small\bfseries] at (Ab) [right=1mm] {$\overline A$};
\node[font=\small\bfseries] at (AB) [right=1mm] {$B$};
\node[font=\small\bfseries] at (ABb) [right=1mm] {$\overline B$};
\node[font=\small\bfseries] at (AbB) [right=1mm] {$B$};
\node[font=\small\bfseries] at (AbBb) [right=1mm] {$\overline B$};
\end{tikzpicture}
\end{center}}



\newcommand{\arbreFicheProbas}{%
\begin{center}
\begin{tikzpicture}[x=0.65cm,y=0.48cm,>=Latex]
\coordinate (O) at (0,0);
\coordinate (A) at (2,1.45);
\coordinate (Ab) at (2,-1.45);
\coordinate (AB) at (4.1,2.15);
\coordinate (ABb) at (4.1,0.75);
\coordinate (AbB) at (4.1,-0.75);
\coordinate (AbBb) at (4.1,-2.15);
\draw[->,line width=.45pt] (O)--(A) node[midway,above left,font=\tiny] {$P(A)$};
\draw[->,line width=.45pt] (O)--(Ab) node[midway,below left,font=\tiny] {$P(\overline A)$};
\draw[->,line width=.45pt] (A)--(AB) node[midway,above,font=\tiny] {$P_A(B)$};
\draw[->,line width=.45pt] (A)--(ABb) node[midway,below,font=\tiny] {$P_A(\overline B)$};
\draw[->,line width=.45pt] (Ab)--(AbB) node[midway,above,font=\tiny] {$P_{\overline A}(B)$};
\draw[->,line width=.45pt] (Ab)--(AbBb) node[midway,below,font=\tiny] {$P_{\overline A}(\overline B)$};
\node[font=\tiny\bfseries,fill=fondpage,inner sep=0.7pt,anchor=south] at ($(A)+(0,0.18)$) {$A$};
\node[font=\tiny\bfseries,fill=fondpage,inner sep=0.7pt,anchor=north] at ($(Ab)+(0,-0.18)$) {$\overline A$};
\node[font=\tiny\bfseries] at (AB) [right=.4mm] {$B$};
\node[font=\tiny\bfseries] at (ABb) [right=.4mm] {$\overline B$};
\node[font=\tiny\bfseries] at (AbB) [right=.4mm] {$B$};
\node[font=\tiny\bfseries] at (AbBb) [right=.4mm] {$\overline B$};
\node[font=\tiny,anchor=west,text=gris] at (4.7,2.15) {$P(A\cap B)$};
\node[font=\tiny,anchor=west,text=gris] at (4.7,0.75) {$P(A\cap\overline B)$};
\node[font=\tiny,anchor=west,text=gris] at (4.7,-0.75) {$P(\overline A\cap B)$};
\node[font=\tiny,anchor=west,text=gris] at (4.7,-2.15) {$P(\overline A\cap\overline B)$};
\end{tikzpicture}
\end{center}
}

\newcommand{\tableauProbas}{%
\begin{center}
\renewcommand{\arraystretch}{1.25}
\begin{tabular}{c|c|c|c}
 & $B$ & $\overline B$ & Total\\ \hline
$A$ & $P(A\cap B)$ & $P(A\cap\overline B)$ & $P(A)$\\ \hline
$\overline A$ & $P(\overline A\cap B)$ & $P(\overline A\cap\overline B)$ & $P(\overline A)$\\ \hline
Total & $P(B)$ & $P(\overline B)$ & $1$
\end{tabular}
\end{center}}

\tcbset{enhanced,arc=2mm,outer arc=2mm,boxrule=0.42pt,left=1.15mm,right=1.15mm,top=0.75mm,bottom=0.75mm,boxsep=0.35mm,before skip=0.8mm,after skip=0.8mm,fonttitle=\bfseries\footnotesize,coltitle=encre,before upper=\raggedright\fontsize{7.65}{8.15}\selectfont}
\newtcolorbox{fichebox}[3][]{colback=white,colframe=#2!72!black,colbacktitle=#2!18,coltitle=#2!50!black,borderline west={0.9mm}{0pt}{#2!78!black},title={#3},drop fuzzy shadow=black!5,#1}
\newcommand{\tagcol}[2]{\begin{tcolorbox}[enhanced,colback=#1!11,colframe=#1!45!black,boxrule=0.42pt,arc=2mm,left=1mm,right=1mm,top=0.45mm,bottom=0.45mm,before skip=0mm,after skip=0.85mm,fontupper=\bfseries\scriptsize,colupper=#1!55!black]\centering #2\end{tcolorbox}}
\newtcolorbox{panel}[1]{colback=fondpage,colframe=fondpage,boxrule=0pt,arc=0mm,left=0mm,right=0mm,top=0mm,bottom=0mm,height=168mm,valign=top,before skip=0mm,after skip=0mm}
\newcommand{\titrepage}{\begin{tcolorbox}[colback=white,colframe=bleu!70!black,boxrule=0.65pt,arc=3mm,left=2.5mm,right=2.5mm,top=1.15mm,bottom=1.15mm,before skip=0mm,after skip=1.4mm,drop fuzzy shadow=black!10]
{\Large\bfseries Parcours de révision -- Probabilités conditionnelles}\hfill {\normalsize Première spécialité -- sans calculatrice}
\end{tcolorbox}}

\begin{document}

\begin{center}
{\LARGE\bfseries Parcours de révision -- Probabilités conditionnelles}\\[1mm]
{\large Première générale -- spécialité mathématiques}\\[1mm]
{\normalsize Sans calculatrice}
\end{center}
\vspace{2mm}
\tableofcontents
\clearpage

\phantomsection
\addcontentsline{toc}{section}{Partie 1 -- Récapitulatif de cours}
\newgeometry{margin=6.5mm}
\pagestyle{empty}
\begin{landscape}
\pagecolor{fondpage}\color{encre}
\thispagestyle{empty}

\fontsize{7.85}{8.45}\selectfont

\titrepage

\noindent
\begin{minipage}[t]{0.242\linewidth}
\tagcol{bleu}{1. EVENEMENTS}
\begin{panel}{bleu}

\begin{fichebox}{bleu}{Vocabulaire}
$\Omega$ est l'univers : ensemble de toutes les issues.

Un évènement est une partie de $\Omega$.
\[
A\cap B : A \text{ et } B
\qquad
A\cup B : A \text{ ou } B.
\]
L'évènement contraire de $A$ est $\overline A$.
\end{fichebox}

\begin{fichebox}{vert}{Formules de base}
\[
P(\overline A)=1-P(A).
\]
Si $A$ et $B$ sont incompatibles :
\[
P(A\cap B)=0.
\]
Dans le cas général :
\[
P(A\cup B)=P(A)+P(B)-P(A\cap B).
\]
\end{fichebox}

\begin{fichebox}{orange}{Partition}
Des évènements forment une \textbf{partition} de l'univers lorsqu'ils sont incompatibles deux à deux et que leur réunion est $\Omega$.

Cas très fréquent :
\[
A \text{ et } \overline A \text{ forment une partition de } \Omega.
\]
\end{fichebox}

\end{panel}
\end{minipage}\hfill
\begin{minipage}[t]{0.242\linewidth}
\tagcol{vert}{2. CONDITIONNEMENT}
\begin{panel}{vert}

\begin{fichebox}{vert}{Probabilité conditionnelle}
Si $P(A)\neq0$, la probabilité de $B$ sachant $A$ est :
\[
P_A(B)=\frac{P(A\cap B)}{P(A)}.
\]
Donc :
\[
P(A\cap B)=P(A)\times P_A(B).
\]
\end{fichebox}

\begin{fichebox}{bleu}{Arbre pondéré}
Sur les branches : probabilités simples ou conditionnelles.

Au bout d'un chemin : probabilité d'une intersection.
\arbreFicheProbas
Exemple : $P(A\cap B)=P(A)\times P_A(B)$.
\end{fichebox}

\begin{fichebox}{rouge}{Attention}
$P_A(B)$ et $P_B(A)$ ne désignent pas la même probabilité.
\[
P_B(A)=\frac{P(A\cap B)}{P(B)}.
\]
\end{fichebox}

\end{panel}
\end{minipage}\hfill
\begin{minipage}[t]{0.242\linewidth}
\tagcol{orange}{3. PROBABILITES TOTALES}
\begin{panel}{orange}

\begin{fichebox}{orange}{Formule}
Si $A$ et $\overline A$ forment une partition de l'univers, alors :
\[
P(B)=P(A\cap B)+P(\overline A\cap B).
\]
Avec un arbre :
\[
P(B)=P(A)P_A(B)+P(\overline A)P_{\overline A}(B).
\]
\end{fichebox}

\begin{fichebox}{violet}{Rédaction attendue}
Toujours écrire :

\emph{$A$ et $\overline A$ forment une partition de l'univers. D'après la formule des probabilités totales, on a :}
\[
P(B)=P(A\cap B)+P(\overline A\cap B).
\]
\end{fichebox}

\begin{fichebox}{bleu}{Tableau}
Un tableau croisé permet de lire directement les intersections et les totaux.
\tableauProbas
\end{fichebox}

\end{panel}
\end{minipage}\hfill
\begin{minipage}[t]{0.242\linewidth}
\tagcol{violet}{4. METHODES}
\begin{panel}{violet}

\begin{fichebox}{violet}{Indépendance}
Deux évènements $A$ et $B$ sont indépendants lorsque :
\[
P(A\cap B)=P(A)P(B).
\]
Si les probabilités conditionnelles existent, c'est équivalent à :
\[
P_A(B)=P(B)
\qquad\text{ou}\qquad
P_B(A)=P(A).
\]
\end{fichebox}

\begin{fichebox}{rouge}{Ne pas confondre}
\textbf{Incompatibles} : ils ne peuvent pas se produire ensemble.
\[
P(A\cap B)=0.
\]
\textbf{Indépendants} : savoir que l'un est réalisé ne change pas la probabilité de l'autre.
\[
P_A(B)=P(B).
\]
\end{fichebox}

\begin{fichebox}{vert}{Choisir la bonne méthode}
\begin{tabularx}{\linewidth}{@{}>{\bfseries}lX@{}}
Arbre & chemin, conditionnement, total\\
Tableau & intersections, marges, pourcentages\\
Question « sachant » & probabilité conditionnelle\\
Question « au total » & formule des probabilités totales\\
Indépendance & comparer $P(A\cap B)$ et $P(A)P(B)$
\end{tabularx}
\end{fichebox}

\end{panel}
\end{minipage}

\end{landscape}
\clearpage
\nopagecolor
\restoregeometry
\pagestyle{fancy}

\section*{Partie 2 -- Questions flash}
\addcontentsline{toc}{section}{Partie 2 -- Questions flash}
\setcounter{question}{0}
\begin{blocquestions}{Questions flash -- Probabilités conditionnelles}
  \item Qu'est-ce que deux évènements incompatibles ? \sourceq{39}
  \item Que signifie $P_B(A)$ ? \sourceq{40}
  \item Quelle est la formule de la probabilité conditionnelle ? \sourceq{41}
  \item Quelle formule permet de retrouver $P(A\cap B)$ à partir de $P(B)$ et $P_B(A)$ ? \sourceq{42}
  \item Que signifie $A\cap B$ ? \sourceq{43}
  \item Que signifie $A\cup B$ ? \sourceq{44}
  \item Que note-t-on souvent pour l'évènement contraire de $A$ ? \sourceq{45}
  \item Que vaut $P(\overline A)$ ? \sourceq{46}
  \item Comment calcule-t-on la probabilité d'un chemin dans un arbre ? \sourceq{47}
  \item Comment calcule-t-on la probabilité d'un évènement obtenu par plusieurs chemins ? \sourceq{48}
  \item Que signifie « $A$ et $B$ forment une partition de l'univers » ? \sourceq{49}
  \item Quelle est la formule des probabilités totales avec une partition $(B_i)$ ? \sourceq{50}
  \item Quand deux évènements $A$ et $B$ sont-ils indépendants ? \sourceq{51}
  \item Si deux évènements sont incompatibles, sont-ils forcément indépendants ? \sourceq{52}
  \item Que signifie un évènement certain ? \sourceq{53}
  \item Quelle est la probabilité de l'évènement impossible ? \sourceq{54}
  \item Si $P(B)\neq0$ et si $P_B(A)=P(A)$, que peut-on conclure ? \sourceq{60}
  \item Si $A$ et $B$ sont incompatibles, que vaut $P(A\cap B)$ ? \sourceq{61}
  \item Qu'appelle-t-on l'univers en probabilités ? \sourceq{64}
  \item Qu'appelle-t-on deux évènements indépendants ? \sourceq{66}
  \item Qu'appelle-t-on une partition de l'univers ? \sourceq{70}
  \item Pour trouver la probabilité d'un évènement $B$ au bout d'un arbre, quelle phrase faut-il écrire avant la formule des probabilités totales ? \sourceq{71}
  \item Dans un exercice, un élève écrit : « $A$ et $B$ sont indépendants, donc $P(A\cup B)=P(A)+P(B)$. » Quelle est l'erreur ? \sourceq{72}
  \item Dans un exercice, un élève écrit : « La probabilité de $A\cap B$, c'est $P(A)\times P(B)$. » Que faut-il préciser ? \sourceq{73}
  \item Si $P(A)=0{,}35$, que vaut $P(\overline A)$ ? \sourceq{74}
  \item Si $A$ et $B$ forment une partition de l'univers, que vaut $P(A)+P(B)$ ? \sourceq{75}
  \item Quels sont les points de rédaction importants pour utiliser la formule des probabilités totales ? \sourceq{76}
  \item Dans un exercice, un élève écrit : « Comme ils n'ont rien à voir, ils sont indépendants. » Comment faut-il justifier l'indépendance ? \sourceq{77}
  \item Après calcul, un élève écrit : « $P_B(A)=P(A)$, donc c'est pareil. » Comment conclure proprement ? \sourceq{78}
  \item Dans un exercice, un élève écrit : « $P(A\cap B)=P(A)+P(B)$. » Pourquoi est-ce faux ? \sourceq{79}
  \item Si $P(A)=\dfrac{3}{5}$ et $P_A(B)=\dfrac{2}{3}$, que vaut $P(A\cap B)$ ?
  \item Si $P(A)=0{,}4$, $P_A(B)=0{,}25$, $P(\overline A)=0{,}6$ et $P_{\overline A}(B)=0{,}5$, que vaut $P(B)$ ?
  \item Si $P(A)=0{,}2$, $P(B)=0{,}5$ et $P(A\cap B)=0{,}1$, les évènements $A$ et $B$ sont-ils indépendants ?
  \item Si $P(B)=0{,}4$ et $P(A\cap B)=0{,}12$, que vaut $P_B(A)$ ?
  \item Dans un arbre, les deux branches issues d'un même noeud portent les probabilités $0{,}7$ et $p$. Que vaut $p$ ?
\end{blocquestions}
\clearpage

\section*{Partie 3 -- QCM d'automatismes}
\addcontentsline{toc}{section}{Partie 3 -- QCM d'automatismes}
\begin{blocquestions}{QCM -- Une seule bonne réponse}
  \item Si $P(A)=0{,}28$, alors $P(\overline A)$ vaut :
  \qcm{$0{,}28$}{$0{,}72$}{$1{,}28$}{$-0{,}28$}
  \item Si $P(A)=\dfrac12$ et $P_A(B)=\dfrac35$, alors $P(A\cap B)$ vaut :
  \qcm{$\dfrac35$}{$\dfrac12$}{$\dfrac3{10}$}{$\dfrac56$}
  \item Si $P(B)=0{,}2$ et $P(A\cap B)=0{,}08$, alors $P_B(A)$ vaut :
  \qcm{$0{,}016$}{$0{,}4$}{$0{,}28$}{$2{,}5$}
  \item Si $A$ et $\overline A$ forment une partition, alors :
  \qcm{$P(A)+P(\overline A)=1$}{$P(A\cap\overline A)=1$}{$P(A)=P(\overline A)$}{$P(A\cup\overline A)=0$}
  \item Si $A$ et $B$ sont indépendants, alors :
  \qcm{$P(A\cup B)=P(A)+P(B)$}{$P(A\cap B)=P(A)P(B)$}{$P(A\cap B)=0$}{$P_A(B)=P(A)$}
  \item Si $P_A(B)=P(B)$ et $P(A)\neq0$, alors :
  \qcm{$A$ et $B$ sont incompatibles}{$A$ et $B$ sont indépendants}{$A=B$}{$P(A)=1$}
\end{blocquestions}
\clearpage


\section*{Partie 4 -- Sujets type bac / E3C}
\addcontentsline{toc}{section}{Partie 4 -- Sujets type bac / E3C}

\titreSujet{Sujet 1 -- Dessert et café}{inspiré de sujets E3C Première sur arbres à trois branches -- sans calculatrice}
Dans un restaurant, les clients choisissent éventuellement un dessert.
\begin{itemize}
  \item $40\,\%$ des clients choisissent un assortiment de macarons ;
  \item $30\,\%$ des clients choisissent une part de tarte ;
  \item les autres clients ne prennent pas de dessert.
\end{itemize}
Aucun client ne prend plusieurs desserts.

Le restaurateur observe que :
\begin{itemize}
  \item parmi les clients qui choisissent des macarons, $75\,\%$ prennent un café ;
  \item parmi les clients qui choisissent une tarte, $\dfrac13$ prennent un café ;
  \item parmi les clients qui ne prennent pas de dessert, $\dfrac16$ prennent un café.
\end{itemize}

On choisit un client au hasard. On note :
\[
\begin{aligned}
M &: \text{« le client choisit des macarons »},\\
T &: \text{« le client choisit une tarte »},\\
N &: \text{« le client ne prend pas de dessert »},\\
C &: \text{« le client prend un café »}.
\end{aligned}
\]
\begin{enumerate}
  \item Justifier que $M$, $T$ et $N$ forment une partition de l'univers.
  \item Construire l'arbre pondéré associé à la situation.
  \item Définir par une phrase les probabilités $P(T\cap C)$ et $P_C(M)$.
  \item Calculer $P(M\cap C)$, $P(T\cap C)$ et $P(N\cap C)$.
  \item En déduire $P(C)$.
  \item On rencontre un client qui a pris un café. Calculer la probabilité qu'il ait choisi des macarons.
  \item Les évènements $M$ et $C$ sont-ils indépendants ? Justifier.
\end{enumerate}

\titreSujet{Sujet 2 -- Atelier et inscription}{inspiré de sujets E3C Première avec tableau croisé -- sans calculatrice}
Lors d'un forum de lycée, $200$ élèves participent à deux types d'activités : un atelier scientifique ou un atelier non scientifique.

On sait que :
\begin{itemize}
  \item $80$ élèves participent à l'atelier scientifique ;
  \item parmi eux, $50$ se sont inscrits à l'avance ;
  \item parmi les élèves qui ne participent pas à l'atelier scientifique, $30$ se sont inscrits à l'avance.
\end{itemize}

On choisit un élève au hasard. On note :
\[
S : \text{« l'élève participe à l'atelier scientifique »},
\qquad
I : \text{« l'élève s'est inscrit à l'avance »}.
\]
\begin{enumerate}
  \item Construire un tableau croisé d'effectifs faisant apparaître $S$, $\overline S$, $I$ et $\overline I$.
  \item Déterminer $P(S)$, $P(I)$ et $P(S\cap I)$.
  \item Calculer $P_S(I)$ et interpréter le résultat.
  \item Calculer $P_I(S)$ et interpréter le résultat.
  \item Comparer $P_S(I)$ et $P(I)$. Que peut-on conjecturer ?
  \item Démontrer si les évènements $S$ et $I$ sont indépendants ou non.
  \item Un élève affirme : « Comme $P_S(I)\neq P_I(S)$, les évènements ne sont pas indépendants. » Expliquer pourquoi cette justification n'est pas correcte.
\end{enumerate}

\titreSujet{Sujet 3 -- Test de contrôle}{inspiré de sujets E3C Première de dépistage / test -- sans calculatrice}
Une entreprise fabrique des pièces mécaniques. On estime que $\dfrac1{20}$ des pièces produites présentent un défaut.

Un test de contrôle est utilisé :
\begin{itemize}
  \item si une pièce est défectueuse, le test est positif dans $\dfrac9{10}$ des cas ;
  \item si une pièce n'est pas défectueuse, le test est positif dans $\dfrac1{10}$ des cas.
\end{itemize}

On choisit une pièce au hasard. On note :
\[
D : \text{« la pièce est défectueuse »},
\qquad
T : \text{« le test est positif »}.
\]
\begin{enumerate}
  \item Donner $P(D)$, $P(\overline D)$, $P_D(T)$ et $P_{\overline D}(T)$.
  \item Construire l'arbre pondéré associé à la situation.
  \item Calculer $P(D\cap T)$.
  \item Calculer $P(\overline D\cap T)$.
  \item Justifier soigneusement que $P(T)=\dfrac7{50}$.
  \item Une pièce a un test positif. Calculer $P_T(D)$.
  \item Interpréter ce résultat dans le contexte.
  \item Un élève affirme : « Le test est positif, donc la pièce est presque certainement défectueuse. » Que peut-on répondre ?
\end{enumerate}
\end{document}
